\documentclass[11pt,a4paper]{article}
\usepackage[margin=1in]{geometry}
\usepackage{amsmath,amssymb,amsthm,booktabs,tabularx,microtype}
\usepackage{hyperref}
\hypersetup{colorlinks=true,linkcolor=black,urlcolor=blue}
\emergencystretch=4em
\newcommand{\R}{\mathbb{R}}
\newcommand{\diag}{\operatorname{diag}}
\newcommand{\blkdiag}{\operatorname{blkdiag}}
\newcommand{\Cov}{\operatorname{Cov}}
\newcommand{\E}{\operatorname{E}}
\newcommand{\Var}{\operatorname{Var}}
\newcommand{\vecrow}{\operatorname{vec}_{\mathrm{row}}}
\setcounter{tocdepth}{1}
\newcolumntype{T}{>{\raggedright\arraybackslash}X}
% Generated by cbmr-proofs. Do not edit.
% Every expression below is rendered from the object sympy verified, not transcribed.
\newcommand{\AppendixPoissonScore}{Y - \mu}
\newcommand{\AppendixPoissonWeight}{\mu}
\newcommand{\AppendixNBScore}{\frac{Y - \mu}{\alpha \mu + 1}}
\newcommand{\AppendixNBObserved}{\frac{\mu \left(Y \alpha + 1\right)}{\left(\alpha \mu + 1\right)^{2}}}
\newcommand{\AppendixNBFisher}{\frac{\mu}{\alpha \mu + 1}}
\newcommand{\AppendixNBWeightDet}{\frac{a \alpha b c d \left(a - b\right) \left(- c + d\right)}{\left(a \alpha c + 1\right) \left(a \alpha d + 1\right) \left(\alpha b c + 1\right) \left(\alpha b d + 1\right)}}

\renewcommand{\thesection}{S\arabic{section}}
\numberwithin{equation}{section}

\title{Supplementary Appendix:\\
Unified Models, Exact Information and Covariance\\
for Spatial Regression and CBMR}
\author{}
\date{}

\begin{document}
\maketitle
\vspace{-2em}
\begin{abstract}
This appendix brings together the lesion-regression manuscript
\texttt{cbmr\_paper.pdf}, the published multi-group CBMR model
\texttt{imag.a.1057.pdf}, and the closed-form Hessian and covariance
notes~\cite{hessian,covariance}.
The central distinction is between CBMR with spatially constant covariate
effects (GC-CBMR) and CBMR with spatially varying covariate effects
(SV-CBMR). These predictor choices are separate from the four observation models.
For each model we specify the likelihood, whether data reduction is exact,
and the resulting information matrix. The shared computational improvements
are weighted basis products, structured inversion, and symmetric
eigenvalue diagnostics. Roughness penalties, shared-covariate uncertainty,
robust covariance and bootstrap inference are treated explicitly rather than
conflated with likelihood simplification. Accompanying SymPy checks verify
exact algebraic identities, not statistical calibration or numerical benchmarks.
\end{abstract}
\tableofcontents
\clearpage

\section{Two CBMR covariate-effect specifications and notation}\label{sec:notation}

\subsection{The key distinction: constant effects, not constant covariate values}
This appendix uses \textbf{GC-CBMR} for CBMR with \emph{globally constant
covariate effects across voxels}, and \textbf{SV-CBMR} for CBMR with
\emph{spatially varying covariate effects}. These are descriptive labels
for the two parameterisations, not different observation distributions.
In both cases the covariate value $Z_{ir}$ may differ between studies.
What changes is the coefficient multiplying it at voxel $j$:
\begin{equation}\label{eq:effect-distinction}
 \frac{\partial\log\mu_{ij}}{\partial Z_{ir}}=
 \begin{cases}
 \gamma_r,&\text{GC-CBMR: constant across }j,\\
 b_r(j)=B_j^\top\beta_r,&\text{SV-CBMR: a spatial coefficient map}.
 \end{cases}
\end{equation}
Here $\beta_r\in\R^P$ is the coefficient block for covariate $r$ within the
stacked SV-CBMR vector $\beta$; group membership and other covariates are
held fixed. For a one-unit change in a continuous covariate, the mean ratio
is $\exp(\gamma_r)$ everywhere under GC-CBMR, but
$\exp\{b_r(j)\}$ under SV-CBMR. The same ratio describes a binary
covariate's change from 0 to 1 when the predictor is linear in that covariate.

\begin{center}\small
\begin{tabularx}{\textwidth}{lTT}
\toprule
Question & GC-CBMR: constant effects & SV-CBMR: varying effects\\
\midrule
What varies spatially? & Group baseline maps $B_j^\top\xi_g$;
the moderator coefficient $\gamma_r$ does not &
The coefficient $b_r(j)$ of each modelled covariate\\
Effect of one covariate & One scalar $\gamma_r$ per specified global effect &
One $P$-coefficient basis expansion $\beta_r$ per effect\\
Interpretation & A common multiplicative scaling of a group's map &
A location-dependent change in the spatial pattern\\
Main computation & Separable mean factors and grouped reductions
(Section~\ref{sec:extensions}) &
Kronecker-design block assembly, generally nonseparable means
(Section~\ref{sec:master})\\
Main covariate test & A global contrast of $\gamma$ &
A voxelwise contrast of $b_r(j)$\\
\bottomrule
\end{tabularx}
\end{center}
\textbf{A spatially varying baseline does not make a covariate effect
spatially varying.} Also, ``constant across voxels'' and ``shared across
groups'' are different restrictions: group-specific coefficients
$\gamma_{gr}$ can still be spatially constant within each group.

\subsection{A matched example with the same groups and one covariate}\label{sec:matched-example}
For a study covariate $x_i$ and the same group baseline
$f_g(j)=B_j^\top\xi_g$, compare
\begin{align}
 \eta_{ij}^{\mathrm{GC}}&=f_{g(i)}(j)+x_i\gamma_x,\\
 \eta_{ij}^{\mathrm{SV}}&=f_{g(i)}(j)+x_i b_x(j),\qquad
 b_x(j)=B_j^\top\beta_x.
\end{align}
GC-CBMR asks whether $x_i$ scales the map up or down by the same ratio
everywhere. SV-CBMR asks \emph{where} its effect is positive, negative or
absent. A positive $\gamma_x$ raises the mean at every voxel, whereas
$b_x(j)$ can have different signs at different locations.

With $G$ groups and a basis of size $P$, these examples have $GP+1$ and
$(G+1)P$ regression coefficients, respectively, before identification
constraints. In the general notation below the GC-CBMR covariate matrix
is $Z=x$ with a separate group-membership matrix $L$; the corresponding
SV-CBMR covariate matrix is $Z=[\,L\ \ x\,]$, whose group-indicator
columns encode the baseline maps. Thus $R$ always means the number of
columns of the \emph{specified} $Z$, not necessarily the same count in
the two designs.

If the basis contains a constant function, $Bc_B=\mathbf1_N$, the
restriction $\beta_x=\gamma_x c_B$ reduces this SV-CBMR example exactly
to GC-CBMR. Constancy is a restriction on the \emph{function}
$B\beta_x$, not in general equality of its basis coefficients. The
corresponding study mean becomes separable; merely observing similar
effect estimates at a few voxels is not an exact separability condition.

\subsection{How the models fit together}
The lesion manuscript~\cite{lesion} and the published CBMR
article~\cite{published} differ in both their \emph{predictor design} and
their \emph{observation likelihood}. These are separate choices.
The former supplies the SV-CBMR specification (Design A below);
the latter supplies the GC-CBMR specification (Design B below).
Poisson, independent NB and clustered NB
likelihoods can be differentiated under either affine predictor.
The efficient aggregated-NB formulas here specifically use the grouped
separable predictor.

\begin{center}\small
\begin{tabularx}{\textwidth}{lTT}
\toprule
Observation model & Data and stochastic assumption & Reduction and destination\\
\midrule
Poisson & Independent cell counts; variance equals mean &
Exact block assembly; grouped model uses both count marginals
(Sections~\ref{sec:families}, \ref{sec:group-poisson})\\
Independent NB & Independent cell counts with variance
$\mu_{ij}+\alpha_g\mu_{ij}^2$ &
Exact cellwise score and weights; generally no marginal-only likelihood
(Section~\ref{sec:families})\\
Aggregated NB & NB approximation to a sum of independent NB cells &
Moment matching is approximate; differentiation of that working likelihood
is exact (Section~\ref{sec:aggregate})\\
Clustered NB & One shared gamma multiplier per study &
Exact marginalisation; correlated voxels and rank-one curvature correction
(Section~\ref{sec:clustered})\\
\bottomrule
\end{tabularx}
\end{center}
Sections~\ref{sec:penalty}--\ref{sec:inference} then apply common estimation
and inference machinery without changing the chosen model.

\subsection{A single notation for both papers}
We retain the principal symbols of Sections 2.2--2.4 of the lesion manuscript.
An observational unit $i$ is a \emph{subject} there and a \emph{study} in CBMR;
$j$ always indexes voxels. The same equations do not imply the same
data-generating assumptions.
\begin{center}
\begin{tabular}{lll}
\toprule
Symbol & Dimension & Meaning\\
\midrule
$M,N$ & scalars & numbers of observational units and voxels\\
$R,P$ & scalars & numbers of covariates and spatial basis functions\\
$Y_{ij},\mu_{ij}$ & scalars & count and mean for unit $i$, voxel $j$\\
$Z,B$ & $M\times R,\ N\times P$ & covariate and spatial design matrices\\
$Z_i^\top,B_j^\top$ & $1\times R,\ 1\times P$ & rows of those matrices\\
$X=Z\otimes B$ & $MN\times PR$ & full design, not explicitly constructed\\
$\beta$ & $PR\times1$ & full-model spatial regression coefficients\\
$\xi,\gamma$ & $P\times1,\ R\times1$ & separable-model spatial and global coefficients\\
$G,\mathcal I_g,M_g$ & group count, set, size & $i\in\mathcal I_g$ iff $g(i)=g$\\
$\xi_g$ & $P\times1$ & group-specific spatial coefficients\\
$\alpha$ & scalar & NB dispersion, $\Var(Y_{ij})=\mu_{ij}+\alpha\mu_{ij}^2$\\
$C$ & $S\times R$ & lesion-model covariate contrast matrix\\
\bottomrule
\end{tabular}
\end{center}
\subsection{SV-CBMR (Design A): spatially varying covariate effects}
Subject $i$ varies slowest in the stacked vectors $Y,\mu$. Define
$\widetilde\beta\in\R^{R\times P}$ by
$\widetilde\beta_{rp}=\beta_{(r-1)P+p}$; thus
$\beta=\vecrow(\widetilde\beta)$. The full predictor is
\begin{equation}\label{eq:full}
 \eta_{ij}=\log\mu_{ij}
 =Z_i^\top\widetilde\beta B_j
 =(Z_i^\top\otimes B_j^\top)\beta,\qquad
 (\eta_{ij})_{i,j}=Z\widetilde\beta B^\top.
\end{equation}
This fixes the vectorisation convention explicitly.
In particular $\beta_r=\widetilde\beta_{r\cdot}^{\top}$ and
$b_r(j)=B_j^\top\beta_r$ in \eqref{eq:effect-distinction}.
An intercept or group-indicator column in $Z$ supplies a spatial
baseline; other columns supply spatially varying moderator effects.

\subsection{GC-CBMR (Design B): spatially constant covariate effects}\label{sec:group-design}
The single-group separable initialisation model in the lesion manuscript is
a \emph{different parameterisation}:
\begin{equation}\label{eq:separable}
 \eta_{ij}=\eta_{Bj}+\eta_{Zi}=B_j^\top\xi+Z_i^\top\gamma,\qquad
 \mu_{Bj}=e^{B_j^\top\xi},\quad \mu_{Zi}=e^{Z_i^\top\gamma},\quad
 \mu_{ij}=\mu_{Bj}\mu_{Zi}.
\end{equation}
The published CBMR model generalises this to
\begin{equation}\label{eq:group-predictor}
 \eta_{ij}=B_j^\top\xi_{g(i)}+Z_i^\top\gamma,\qquad
 \mu_{ij}=\mu_{B,g(i),j}\mu_{Zi},\qquad
 \mu_{Bgj}=\exp(B_j^\top\xi_g).
\end{equation}
Let $L_{ig}=\mathbf1\{g(i)=g\}$ be the $M\times G$ membership matrix,
$\xi_{\mathrm{all}}=(\xi_1^\top,\ldots,\xi_G^\top)^\top$, and
$\vartheta=(\xi_{\mathrm{all}}^\top,\gamma^\top)^\top$. Then
\begin{equation}\label{eq:group-design}
 \eta=\mathcal X\vartheta,\qquad
 \mathcal X=[\,L\otimes B\ \ Z\otimes\mathbf1_N\,],
 \qquad \dim(\vartheta)=GP+R.
\end{equation}
This restates Equations 1--5 of the published article. Set $G=1$ to obtain
\eqref{eq:separable}; omit the global block when there are no moderators.
Group-specific global effects may be encoded by group-supported columns of
$Z$; in that case $R$ counts all such columns. Separate group fits instead
have separate parameter vectors and independent estimates when their data
are independent.

We use $\ell(\beta,\alpha)$ for the full-model log-likelihood and
$\ell(\xi,\gamma,\alpha)$ for the separable model; $\gamma$ is not substituted
for the full coefficient vector $\beta$.

\subsection{Translation from the published article and source notes}
\begin{center}\small
\begin{tabularx}{\textwidth}{llT}
\toprule
Published CBMR symbol & Symbol here & Reason\\
\midrule
$X,\ x_j^\top$ & $B,\ B_j^\top$ & Reserve $X=Z\otimes B$ for Design A\\
$\beta_g$ & $\xi_g$ & Reserve $\beta$ for spatially varying covariate effects\\
$\mu_g^X,\ \mu^Z$ & $\mu_{Bg},\ \mu_Z$ & Same spatial/global mean factors\\
$Y_{ti},\ Y_{gj}$ & $Y_{i\cdot},\ Y_{gj}$ & Study total and group--voxel total\\
$r'_{gj},\ p'_{gj},\ \alpha'_g$ & $\rho_g,\ p_{gj},\ \alpha_g^{\mathrm{eff}}$ &
Aggregated NB shape, probability and effective dispersion\\
$\lambda_i$ & $\Lambda_i$ & Random study multiplier, not a smoothing weight\\
$C\ (m\times S)$ & $C_G\ (S\times G)$ &
$S$ always counts tested constraints here\\
$J$ & $J_B$ & Roughness matrix; distinct from sandwich meat $\mathcal J$\\
\bottomrule
\end{tabularx}
\end{center}
In the source CBMR notes, $(n,V,K)$ correspond to $(M,N,P)$ here. Their
non-spatial design $G$ becomes $Z$ in \eqref{eq:separable}, and their
$(u_v,e_i,\theta)$ become $(\mu_{Bj},\mu_{Zi},\alpha)$. Pattern counts are
denoted by $G$ for grouped CBMR, not by $P$, which is reserved
for basis size. Their NB shape $R_p$ is renamed $\rho_g$ to avoid collision
with the lesion manuscript's covariate count $R$. $C$ remains a contrast matrix and is
not used for an information cross block. $C_G$ and $C_\gamma$ distinguish
group and global-covariate contrasts from the lesion manuscript's $C$.

\subsection{Scope and inferential assumptions}
All regression derivatives hold dispersion fixed, including when it is set to
an estimate $\widehat\alpha$ (or $\widehat\alpha_g$). They do not propagate dispersion-estimation
uncertainty. If joint observed-information inference is required, the
dispersion row and column must also be included before inversion. All inverses
below require an identifiable parameterisation and nonsingular stated blocks.
Mean-centering covariates does not by itself establish full rank, especially
when covariates are constant within groups or a basis has redundant columns.
Binary data do not make a Poisson or NB working likelihood identical to a
Bernoulli likelihood.

\section{SV-CBMR: exact design and Hessian assembly}\label{sec:master}

For SV-CBMR, let the negative log-likelihood be
$F(\beta)=\sum_{i,j}f_{ij}(\eta_{ij})$, holding all non-regression quantities
fixed. Write $s_{ij}=-f'_{ij}(\eta_{ij})$ and
$w_{ij}=f''_{ij}(\eta_{ij})$. With
$X_i=Z_i^\top\otimes B\in\R^{N\times PR}$ and
$W_i^{\mathrm{obs}}=\diag(w_{i1},\ldots,w_{iN})$,
\begin{align}
 \nabla_\beta\ell &= \sum_i X_i^\top s_i, \label{eq:score}\\
 H=-\nabla_\beta^2\ell
 &=\sum_i X_i^\top W_i^{\mathrm{obs}}X_i
 =\sum_i (Z_iZ_i^\top)\otimes
       \bigl(B^\top W_i^{\mathrm{obs}}B\bigr), \label{eq:master}\\
 H_{(r,p),(r',p')}
 &=\sum_i Z_{ir}Z_{ir'}\sum_j w_{ij}B_{jp}B_{jp'}.
 \label{eq:entry}
\end{align}
\begin{proof}
The derivative of $\eta_{ij}$ in $\beta_{(r,p)}$ is $Z_{ir}B_{jp}$
and every second derivative of $\eta_{ij}$ is zero. Applying the chain rule
to each $f_{ij}$ therefore gives \eqref{eq:entry}, with no predictor-curvature
term. The Kronecker mixed-product identity gives \eqref{eq:master}.
\end{proof}
This is the source notes' affine-predictor lemma in the unified notation.
No independence claim about real lesion data is needed to differentiate the
working likelihood. Its statistical consequences are addressed by the sandwich
estimator in Section~\ref{sec:inference}.

Equation~\eqref{eq:master} can be evaluated subjectwise or in batches. It does
not require $X$, an $MN\times MN$ diagonal matrix, or a
$(PR)\times M\times N$ derivative intermediate. A weighted basis product
$B^\top\diag(w_i)B$ costs $O(NP^2)$ for a dense basis, with an $N\times P$
workspace; scatter into the $R^2$ coefficient blocks costs $O(R^2P^2)$
per subject. Dense information storage still costs $O((PR)^2)$.
If many subjects have exactly the same row of $Z$, their weights may first be
summed at each voxel and the weighted basis product computed once per distinct
row. No approximation is involved.
The affine chain rule is shared by both specifications.
For GC-CBMR (Design B) use the corresponding study rows $\mathcal X_i$ from
\eqref{eq:group-design}. More generally, if the study likelihood depends
jointly on its log-mean vector, the same affine chain rule gives
$H=\sum_i\mathcal X_i^\top W_i^{\mathrm{obs}}\mathcal X_i$ with
$W_i^{\mathrm{obs}}=\nabla_{\eta_i}^2(-\ell_i)$.
Independent-cell models make $W_i^{\mathrm{obs}}$ diagonal; clustered NB
has the diagonal-minus-rank-one form in
\eqref{eq:cluster-local-information}.

\section{Observation likelihoods shared by both specifications}\label{sec:families}

The Poisson or NB choice does not determine whether the covariate effect
is spatially constant or varying. The scalar derivatives below apply to
either predictor. Projection into regression coefficients uses $X_i$
for SV-CBMR and $\mathcal X_i$ for GC-CBMR.

\subsection{Scalar derivatives and Fisher weights}
For one cell, suppress indices and put $\mu=e^\eta$. With parameter-free
terms retained or discarded consistently, the two log-likelihoods are
\begin{align}
 \ell_{\mathrm{P}}&=Y\eta-e^\eta-\log\Gamma(Y+1),\\
 \ell_{\mathrm{NB}}&=\log\Gamma(Y+\alpha^{-1})-\log\Gamma(\alpha^{-1})
 -\log\Gamma(Y+1)+Y\log\alpha+Y\eta
 -(Y+\alpha^{-1})\log(1+\alpha e^\eta).
 \label{eq:nb-ll}
\end{align}
For $Y\in\{0,1\}$, the gamma terms and $Y\log\alpha$ cancel, recovering
Equation 2.10 of the lesion manuscript. The exact scalar derivatives are
\begin{align}
 s_{\mathrm{P}}=\frac{\partial\ell_{\mathrm{P}}}{\partial\eta}
   &=\AppendixPoissonScore,&
 w_{\mathrm{P}}=-\frac{\partial^2\ell_{\mathrm{P}}}{\partial\eta^2}
   &=\AppendixPoissonWeight,\label{eq:poisson-weights}\\
 s_{\mathrm{NB}}=\frac{\partial\ell_{\mathrm{NB}}}{\partial\eta}
   &=\AppendixNBScore,&
 w_{\mathrm{NB}}^{\mathrm{obs}}
   =-\frac{\partial^2\ell_{\mathrm{NB}}}{\partial\eta^2}
   &=\AppendixNBObserved,\label{eq:nb-weights}\\
 w_{\mathrm{NB}}^{\mathrm{F}}
   =\E[w_{\mathrm{NB}}^{\mathrm{obs}}]
   &=\AppendixNBFisher.&&\label{eq:fisher-weight}
\end{align}
The right-hand sides of \eqref{eq:poisson-weights}--\eqref{eq:fisher-weight}
are generated from the SymPy objects checked by \texttt{proofs/appendix.py}.
The expectation uses only $\E[Y]=\mu$, since the observed NB weight is affine
in $Y$. It agrees with
$\Var(s_{\mathrm{NB}})=(\mu+\alpha\mu^2)/(1+\alpha\mu)^2$ under the NB model.
Both NB weights and the NB score have the Poisson limits as $\alpha\to0^+$.

Thus \eqref{eq:master} with $w^{\mathrm{obs}}$ gives the observed information
$H$. Replacing it by $w^{\mathrm F}$ gives the Fisher information
\begin{align}\label{eq:fisher}
 I(\beta)=\sum_i X_i^\top W_iX_i,\qquad
 W_i&=\diag(W_{i1},\ldots,W_{iN}),\\
 W_{ij}&=
 \begin{cases}
 \mu_{ij},&\text{Poisson},\\
 \mu_{ij}/(1+\alpha\mu_{ij}),&\text{independent NB}.
 \end{cases}\notag
\end{align}
These match the weights in the lesion manuscript's Equation 2.17.
For grouped independent NB use $\alpha_{g(i)}$ in the same formulas.
For Poisson $H=I$ and
neither contains $Y$ at fixed $\beta$; evaluated at $\widehat\beta$, both
depend indirectly on the data through the fit. For NB, $H$ and $I$ generally
differ even at a finite-sample optimum.

\subsection{The NB score is not the Poisson residual score}
The same log link does \emph{not} give the same score for the two families:
\begin{equation}\label{eq:nb-score}
 U_i=\nabla_\beta\ell_i=
 \begin{cases}
 X_i^\top(Y_i-\mu_i),&\text{Poisson},\\
 X_i^\top\diag\bigl((1+\alpha\mu_i)^{-1}\bigr)(Y_i-\mu_i),
       &\text{independent NB}.
 \end{cases}
\end{equation}
Vector divisions and powers here and below are elementwise.
For NB Fisher scoring the exact update direction is
$I(\beta)^{-1}\sum_i U_i$. Equivalently, the IRLS working response is
$\eta+(Y-\mu)/\mu$ and multiplying its residual by the Fisher weights gives
$(Y-\mu)/(1+\alpha\mu)$, not $Y-\mu$.
Consequently the unweighted NB residual in the lesion manuscript's update and
score discussion must not be used as an NB likelihood score. This appendix
retains the manuscript's notation while making that mathematical distinction explicit.

\subsection{What separability does and does not simplify}
At a mean of the form $\mu_{ij}=\mu_{Zi}\mu_{Bj}$, the \emph{full-model}
Poisson information has the exact factorisation
\begin{equation}\label{eq:kron-information}
 I_{\mathrm P}(\beta)
 =(Z^\top W_ZZ)\otimes(B^\top W_BB),\qquad
 W_Z=\diag(\mu_Z),\quad W_B=\diag(\mu_B).
\end{equation}
This follows by applying the mixed-product identity to
$X^\top(W_Z\otimes W_B)X$.
When both factors are nonsingular its inverse is the Kronecker product of
their inverses. Factorisation costs $O(R^3+P^3)$ instead of $O((PR)^3)$;
materialising a dense covariance still requires $O((PR)^2)$ storage.
Fixing these weights at an initial fit gives a preconditioner, not the
current information of a nonseparable full-model fit. Convergence additionally
requires a suitable optimisation procedure; a fixed preconditioner alone
does not guarantee convergence.

For NB, even separable means do not generally yield separable Fisher weights.
For two subjects with factors $a,b$ and two voxels with factors $c,d$, let
$W^*_{ij}=\mu_{Zi}\mu_{Bj}/(1+\alpha\mu_{Zi}\mu_{Bj})$. Then
\begin{equation}\label{eq:nonseparable}
 \det W^*=\AppendixNBWeightDet.
\end{equation}
This is generically nonzero for positive factors, $\alpha>0$, $a\ne b$ and
$c\ne d$. A separable weight grid has rank one, so the Poisson factorisation
cannot simply be relabelled as an exact NB factorisation. The block assembly
\eqref{eq:master} remains exact for either NB weight.

\section{GC-CBMR: grouped reductions and information}\label{sec:extensions}

\subsection{Common grouped design and count summaries}
Use GC-CBMR (Design B) throughout this section. Its spatially constant
covariate effects yield the factorisation
$\mu_{ij}=\mu_{B,g(i),j}\mu_{Zi}$.
The grouped Poisson compression and aggregated-NB reparameterisation
below rely on this structure and must not be transferred unchanged to
unrestricted SV-CBMR. The local clustered-NB weights
\eqref{eq:cluster-local-information}, by contrast, apply to either
affine predictor when projected with the appropriate design.
Define
$Y_{gj}=\sum_{i\in\mathcal I_g}Y_{ij}$,
$Y_{i\cdot}=\sum_jY_{ij}$, and
$Y_{g\cdot}=\sum_{i\in\mathcal I_g}Y_{i\cdot}=\sum_jY_{gj}$.
The vector $Y_g$ contains the group--voxel counts, not the full stacked
observations. Set $Z_g$ to the rows of $Z$ in $\mathcal I_g$ and
$\mu_{Zg}$ to the corresponding global mean factors.

For one group, suppress $g$ and restrict study sums to $\mathcal I_g$.
The cell design row for its spatial and shared global coefficients is
$(B_j^\top,Z_i^\top)$.
For any cellwise likelihood with observed weight $w_{ij}$,
\begin{align}
 H_{\xi\xi}&=B^\top\diag\left(\left\{\sum_i w_{ij}\right\}_j\right)B,\\
 H_{\xi\gamma}&=\sum_{i,j}w_{ij}B_jZ_i^\top,\\
 H_{\gamma\gamma}&=Z^\top\diag\left(\left\{\sum_j w_{ij}\right\}_i\right)Z.
 \label{eq:separable-blocks}
\end{align}
Use \eqref{eq:poisson-weights} or \eqref{eq:nb-weights} for the observed
information and \eqref{eq:fisher-weight} for independent NB Fisher information.
These are $(P+R)\times(P+R)$ information matrices, not the $PR\times PR$
matrix in \eqref{eq:kron-information}. In the joint grouped fit, place each
spatial block in its own position, retain its cross block with $\gamma$,
and sum the global blocks over groups. Clustered NB is not cellwise
independent and instead uses Section~\ref{sec:clustered}.

\subsection{Poisson: exact marginal compression, not loss of study totals}\label{sec:group-poisson}
\paragraph{Likelihood and exact reduction.}
Define $E_g=\sum_j\mu_{Bgj}$, $T_g=\sum_{i\in\mathcal I_g}\mu_{Zi}$,
$a_g=B^\top\mu_{Bg}$ and $u_g=Z_g^\top\mu_{Zg}$.
Substituting the grouped predictor into the independent Poisson
log-likelihood and collecting terms gives, up to data-only constants,
\begin{equation}\label{eq:group-poisson}
 \ell_{\mathrm P}
 =\sum_gY_g^\top\log\mu_{Bg}
   +\sum_iY_{i\cdot}\log\mu_{Zi}-\sum_g E_gT_g.
\end{equation}
This is the computationally useful marginal expression in Equation 6 of
the published article: both $Y_g$ and $Y_{i\cdot}$ must be retained when
global covariates are estimated. Computing the mean term and scores uses
group--voxel arrays and study vectors rather than an $M\times N$ mean array.

There is an important distinction between this compression and fitting only
$Y_{gj}\sim\operatorname{Poisson}(\mu_{Bgj}T_g)$. Up to data-only constants,
\begin{equation}\label{eq:poisson-allocation}
 \ell_{\mathrm{full}}-\ell_{\mathrm{totals}}
 =\sum_g\sum_{i\in\mathcal I_g}Y_{i\cdot}
                   \log\left(\frac{\mu_{Zi}}{T_g}\right).
\end{equation}
The missing term is the conditional multinomial allocation of each voxel's
total across studies, with probabilities $\mu_{Zi}/T_g$. Its global score is
\[
 \sum_g\left\{\sum_{i\in\mathcal I_g}Y_{i\cdot}Z_i
                         -Y_{g\cdot}\frac{u_g}{T_g}\right\}.
\]
Thus voxel totals alone do not generally give the same covariate likelihood.
If the allocation probabilities are fixed rather than estimated, that term
is parameter-free for the spatial fit.

\paragraph{Exact information and scores.}
The nonzero blocks reduce to
\begin{align}\label{eq:separable-poisson}
 H_{\xi_g\xi_g}&=T_g B^\top\diag(\mu_{Bg})B,&
 H_{\xi_g\gamma}&=a_gu_g^\top,\\
 H_{\gamma\gamma}&=\sum_gE_g Z_g^\top\diag(\mu_{Zg})Z_g.\notag
\end{align}
The scores are
\[
 \nabla_{\xi_g}\ell=B^\top(Y_g-T_g\mu_{Bg}),\qquad
 \nabla_\gamma\ell=\sum_iZ_i(Y_{i\cdot}-E_{g(i)}\mu_{Zi}).
\]
For $G=1$ these are the separable-initialisation formulas of the lesion
manuscript. Differentiation is exact; faster evaluation does not change
the Poisson approximation used for binary data.

\subsection{Aggregated NB: moment-matched totals and exact working curvature}\label{sec:aggregate}
\paragraph{Likelihood and approximation boundary.}
The NB calculation in \texttt{cbmr\_hessian.pdf} is not
\eqref{eq:nb-ll} summed over individual subjects and voxels.
It is the moment-matched NB likelihood for group--voxel totals in
Equations 7--9 of the published article. Define
\begin{equation}
 s_{1g}=\sum_{i\in\mathcal I_g}\mu_{Zi},\qquad
 s_{2g}=\sum_{i\in\mathcal I_g}\mu_{Zi}^2,\qquad
 A_g=\frac{s_{1g}}{\alpha_gs_{2g}},\qquad
 \rho_g=\frac{s_{1g}^2}{\alpha_gs_{2g}}=s_{1g}A_g.
\end{equation}
Here $\alpha_g>0$ is a group-specific extension of the lesion manuscript's global
$\alpha$; setting all $\alpha_g=\alpha$ recovers a shared dispersion.
The summed independent NB counts have moments
\begin{align}
 m_{gj}=\E[Y_{gj}]&=\mu_{Bgj}s_{1g},\\
 \Var(Y_{gj})&=\mu_{Bgj}s_{1g}
                    +\alpha_g\mu_{Bgj}^2s_{2g}.
\end{align}
Using the NB convention
$\Pr(Y=y)=\Gamma(y+\rho)/[\Gamma(\rho)\Gamma(y+1)](1-p)^\rho p^y$,
the matching distribution has
\begin{equation}\label{eq:effective}
 p_{gj}=\frac{\mu_{Bgj}}{\mu_{Bgj}+A_g},\qquad
 \alpha_g^{\mathrm{eff}}=\frac{\alpha_gs_{2g}}{s_{1g}^2}
 =\frac1{\rho_g},\qquad
 \Var_{\mathrm{NB}}(Y_{gj})=m_{gj}
                       +\alpha_g^{\mathrm{eff}}m_{gj}^2.
\end{equation}
This reduces the article's apparently voxel-specific shape $r'_{gj}$
to a single shape $\rho_g$ per group. The spatial factor cancels exactly.
It does not turn an arbitrary NB convolution into an NB distribution.
For example, writing $s_{3g}=\sum_{i\in\mathcal I_g}\mu_{Zi}^3$,
the third cumulant of the true sum minus that of the matched NB is
\begin{equation}
 2\alpha_g^2\mu_{Bgj}^3
       \left(s_{3g}-\frac{s_{2g}^2}{s_{1g}}\right).
\end{equation}
If every study factor in a group is equal, the independent NB summands
have a common success probability and the NB convolution is exact;
otherwise the higher cumulants need not match.

Suppress $g$ within the following formulas and write
$d_j=\mu_{Bj}+A$. Up to a data-only constant, the negative log-likelihood is
\begin{equation}\label{eq:psi}
 \begin{split}
 \Psi(\rho,A,\eta_B)={}&-\sum_j\log\Gamma(Y_{gj}+\rho)
 +N\log\Gamma(\rho)-N\rho\log A\\
 &+\sum_j(\rho+Y_{gj})\log d_j-\sum_jY_{gj}\eta_{Bj}.
 \end{split}
\end{equation}
\paragraph{Exact information for the matched likelihood.}
It has diagonal log-spatial curvature
\begin{equation}\label{eq:aggregate-weight}
 w_j=\frac{(\rho+Y_{gj})\mu_{Bj}A}{d_j^2},\qquad
 H_{\xi_g\xi_g}=B^\top\diag(w)B.
\end{equation}
For global coefficients, let $\rho_r=\partial\rho/\partial\gamma_r$ and
$A_r=\partial A/\partial\gamma_r$, with analogous second derivatives.
Then
\begin{align}
 (H_{\xi_g\gamma})_{pr}
 &=\sum_j B_{jp}\left\{\frac{\mu_{Bj}}{d_j}\rho_r
       -\frac{(\rho+Y_{gj})\mu_{Bj}}{d_j^2}A_r\right\},\\
 (H_{\gamma\gamma}^{(g)})_{rr'}
 &=\Psi_{\rho\rho}\rho_r\rho_{r'}
 +\Psi_{\rho A}(\rho_rA_{r'}+A_r\rho_{r'})
 +\Psi_{AA}A_rA_{r'}+\Psi_\rho\rho_{rr'}+\Psi_A A_{rr'}.
 \label{eq:aggregate-chain}
\end{align}
The global block is summed over groups. With $\psi,\psi_1$ denoting digamma
and trigamma,
\begin{align}
 \Psi_\rho&=-\sum_j\psi(Y_{gj}+\rho)+N\psi(\rho)-N\log A+\sum_j\log d_j,\\
 \Psi_A&=-N\rho/A+\sum_j(\rho+Y_{gj})/d_j,\\
 \Psi_{\rho\rho}&=-\sum_j\psi_1(Y_{gj}+\rho)+N\psi_1(\rho),\\
 \Psi_{\rho A}&=-N/A+\sum_jd_j^{-1},\qquad
 \Psi_{AA}=N\rho/A^2-\sum_j(\rho+Y_{gj})/d_j^2.
\end{align}
For an explicit evaluation of the inner-map derivatives, put
\[
 u_1=\sum_i\mu_{Zi}Z_i,\quad u_2=\sum_i\mu_{Zi}^2Z_i,\quad
 \mathcal M_1=\sum_i\mu_{Zi}Z_iZ_i^\top,\quad
 \mathcal M_2=\sum_i\mu_{Zi}^2Z_iZ_i^\top,
\]
with sums over $\mathcal I_g$, and define
\[
 l=u_1/s_1,\quad h=2u_2/s_2,\quad
 \mathcal L_1=\mathcal M_1/s_1-u_1u_1^\top/s_1^2,\quad
 \mathcal L_2=4\mathcal M_2/s_2-4u_2u_2^\top/s_2^2.
\]
Then
\begin{align}
 \nabla A&=A(l-h),&
 \nabla^2 A&=A\{(l-h)(l-h)^\top+\mathcal L_1-\mathcal L_2\},\\
 \nabla\rho&=\rho(2l-h),&
 \nabla^2\rho&=\rho\{(2l-h)(2l-h)^\top+2\mathcal L_1-\mathcal L_2\}.
\end{align}
These follow by differentiating
$\log A=\log s_1-\log s_2-\log\alpha_g$ and
$\log\rho=2\log s_1-\log s_2-\log\alpha_g$.
They retain the inner-map curvature terms in \eqref{eq:aggregate-chain}.
The spatial weight tends to $s_1\mu_{Bj}$ as $\alpha_g\to0^+$.
The reparameterisation is exact for this aggregated likelihood; it does not
make the moment-matching approximation exact for arbitrary independent NB sums.
The derivatives hold the \emph{original} $\alpha_g$ fixed, not
$\alpha_g^{\mathrm{eff}}$. If instead effective dispersion is independently
parameterised and held fixed, $\rho_g$ is constant and $A_g=\rho_g/s_{1g}$;
its derivatives must be recomputed for that convention. Mixing these two
conditioning choices changes the Hessian.

\subsection{Clustered NB: exact study marginalisation and rank-one curvature}\label{sec:clustered}
\paragraph{Likelihood and dependence.}
The clustered extension places a shared gamma effect within each subject:
$\Lambda_i\sim\operatorname{Gamma}(\nu_g,\text{rate }\nu_g)$,
$\nu_g=1/\alpha_g$, and, conditionally on $\Lambda_i$,
$Y_{ij}$ are independent Poisson with mean $\Lambda_i\mu_{Bgj}\mu_{Zi}$ for
$i\in\mathcal I_g$.
Independent $\Lambda_i$ across studies give
\begin{equation}\label{eq:cluster-covariance}
 \Cov(Y_i)=\diag(\mu_i)+\alpha_g\mu_i\mu_i^\top,\qquad
 \Cov(Y_i,Y_{i'})=0\quad(i\ne i').
\end{equation}
This includes both the diagonal variance and the off-diagonal covariance
in Equation 10 of the published article; the latter is absent under
independent NB.

Let $t_i=\sum_j\mu_{ij}=E_g\mu_{Zi}$. Exact marginalisation can be understood
as a total-count model followed by spatial allocation:
\begin{equation}\label{eq:cluster-factorisation}
 Y_{i\cdot}\sim\mathrm{NB}(\text{mean }t_i,\text{shape }\nu_g),\qquad
 Y_i\mid Y_{i\cdot}\sim
 \mathrm{Multinomial}\left(Y_{i\cdot},
                   \left\{\frac{\mu_{Bgj}}{E_g}\right\}_j\right).
\end{equation}
Multiplying these two mass functions cancels the total factorial and the
allocation powers of $t_i$, giving the joint gamma--Poisson mass function.
Unlike aggregated NB, this reduction is exact. Collecting its terms over
studies gives Equation 11 of the article, up to data-only terms:
\begin{equation}
 \begin{split}
 \ell_g={}&|\mathcal I_g|\nu_g\log\nu_g-|\mathcal I_g|\log\Gamma(\nu_g)
 +\sum_{i\in\mathcal I_g}\log\Gamma(Y_{i\cdot}+\nu_g)\\
 &-\sum_{i\in\mathcal I_g}(Y_{i\cdot}+\nu_g)\log(E_g\mu_{Zi}+\nu_g)
 +\sum_jY_{gj}\eta_{Bgj}+\sum_{i\in\mathcal I_g}Y_{i\cdot}\eta_{Zi},
 \end{split}
\end{equation}
where $E_g=\sum_j\mu_{Bgj}$. Write
$d_i=E_g\mu_{Zi}+\nu_g$, $a_g=B^\top\mu_{Bg}$ and
\[
 g'_g=\sum_{i\in\mathcal I_g}\frac{(Y_{i\cdot}+\nu_g)\mu_{Zi}}{d_i},
 \qquad
 g''_g=-\sum_{i\in\mathcal I_g}\frac{(Y_{i\cdot}+\nu_g)\mu_{Zi}^2}{d_i^2}.
\]
\paragraph{Exact observed information.}
Then
\begin{align}
 \nabla_{\eta_{Bg}}^2(-\ell_g)
 &=g'_g\diag(\mu_{Bg})+g''_g\mu_{Bg}\mu_{Bg}^\top,\\
 H_{\xi_g\xi_g}
 &=g'_g B^\top\diag(\mu_{Bg})B+g''_g a_ga_g^\top,\\
 H_{\xi_g\gamma}
 &=a_g\left\{\sum_{i\in\mathcal I_g}
       \frac{(Y_{i\cdot}+\nu_g)\nu_g\mu_{Zi}}{d_i^2}Z_i^\top\right\},\\
 H_{\gamma\gamma}^{(g)}
 &=\sum_{i\in\mathcal I_g}
       \frac{(Y_{i\cdot}+\nu_g)\nu_g E_g\mu_{Zi}}{d_i^2}Z_iZ_i^\top.
 \label{eq:clustered}
\end{align}
The rank-one term is a negative correction since $g''_g<0$.
Only subject totals enter these Hessian weights: the terms containing the
voxel marginals are linear in the coefficients and differentiate away.
The result follows from differentiating a scalar function of
$E_g=\sum_j e^{\eta_{Bgj}}$ twice. Thus the basis projection still needs only
a weighted Gram matrix and one outer product, not an $N\times N$ matrix.

\paragraph{Expected information and fixed-dispersion curvature.}
For a single study expressed in its full log-mean vector $\eta_i$,
the observed information before basis projection is
\begin{equation}\label{eq:cluster-local-information}
 W_i^{\mathrm{obs}}
 =\frac{Y_{i\cdot}+\nu_g}{t_i+\nu_g}
       \left\{\diag(\mu_i)-\frac{\mu_i\mu_i^\top}{t_i+\nu_g}\right\},
 \qquad
 W_i^{\mathrm F}=\diag(\mu_i)-\frac{\mu_i\mu_i^\top}{t_i+\nu_g}.
\end{equation}
The expectation uses $\E[Y_{i\cdot}]=t_i$. These weights are \emph{matrices},
not the diagonal independent-NB weights. For any real vector $v$,
\begin{equation}\label{eq:cluster-positive}
 (t_i+\nu_g)v^\top W_i^{\mathrm F}v
 =\nu_g\sum_j\mu_{ij}v_j^2
   +\sum_{j<k}\mu_{ij}\mu_{ik}(v_j-v_k)^2.
\end{equation}
Consequently the conditional regression objective is convex for fixed
$\nu_g>0$ and an affine log predictor. The same is true for independent
Poisson and NB with fixed dispersion. This is not a claim of joint convexity
when dispersion is also estimated, nor of convexity of the
moment-matched NB composition in $(\xi_g,\gamma)$.

\section{Roughness penalties and the fitted objective}\label{sec:penalty}

\subsection{An explicit quadratic-penalty convention}
The published article introduces a spatial roughness penalty in Section 2.1.1
and Figure 2. To carry that improvement through the Hessian calculation,
let $J_B=J_B^\top\succeq0$ be a fixed $P\times P$ roughness matrix and
$\lambda_g\geq0$ fixed smoothing weights. We use the explicit convention
\begin{equation}
 \mathcal P(\vartheta)=\frac12\sum_g\lambda_g\xi_g^\top J_B\xi_g
 =\frac12\vartheta^\top K\vartheta,\qquad
 K=\blkdiag(\lambda_1J_B,\ldots,\lambda_GJ_B,0_{R\times R}).
\end{equation}
The factor $1/2$ is a convention: a penalty written without it has twice the
gradient and Hessian for the same numerical $\lambda_g$. No particular
roughness operator or smoothing weight is inferred from the PDF; its detailed
construction is referred to Supplementary Appendix A.1 of the article,
which is not included in the supplied PDF.

For the minimised objective $F_{\mathrm{pen}}=-\ell+\mathcal P$,
\begin{equation}\label{eq:penalty}
 \nabla F_{\mathrm{pen}}=-\nabla\ell+K\vartheta,\qquad
 H_{\mathrm{pen}}=H+K.
\end{equation}
Thus the closed-form information receives a fixed block-diagonal addition;
the shared-global border and cross-group structural zeros are preserved.
For SV-CBMR (Design A) use
$K=\blkdiag(\lambda_1J_B,\ldots,\lambda_RJ_B)$ on the $R$ spatial maps.
A roughness penalty need not penalise constant directions and therefore
need not remove an intercept-identification problem.

\subsection{Curvature is not automatically a sampling covariance}
For an unpenalised likelihood, model-based covariance uses inverse
information. For a penalised estimator, inverse penalised curvature is
not generally its frequentist sampling covariance. Linearising the
penalised estimating equation gives
\begin{equation}\label{eq:penalty-covariance}
 V_{\mathrm{pen}}\ \approx\
 (H+K)^{-1}\mathcal J(H+K)^{-\top},
\end{equation}
where $\mathcal J$ is the variance estimate of the \emph{unpenalised}
data score, not of a score with the penalty repeatedly added per study.
Under a correctly specified model replace the bread $H$ and score variance
by $I$ to obtain $(I+K)^{-1}I(I+K)^{-1}$. In particular,
\begin{equation}
 (I+K)^{-1}-(I+K)^{-1}I(I+K)^{-1}
 =(I+K)^{-1}K(I+K)^{-1}.
\end{equation}
This identity distinguishes penalised curvature from sampling variation;
neither expression by itself accounts for smoothing bias or data-driven
smoothing-parameter uncertainty.
Alternating dispersion and regression updates, as in the article's
Algorithm 1, does not remove these distinctions. A final inference step
must state whether it uses likelihood information, penalised curvature,
or a penalised estimating-equation sandwich.

\section{From information to covariance}\label{sec:covariance}

The following linear-algebra identities apply to a symmetric nonsingular
information matrix, whether $H$ is observed information or $I$ is the Fisher
information. They do not change which statistical estimator has been chosen.

\subsection{Structural zeros are determined by the design}
For the full independent-cell model, \eqref{eq:entry} shows that the
$(r,r')$ coefficient block vanishes if no subject loads both covariates:
$Z_{ir}Z_{ir'}=0$ for every $i$. More generally, connect two covariate
columns when they co-occur in any row of $Z$, and take connected components.
Different components have zero cross blocks, so after reordering,
\begin{equation}
 H=\blkdiag(H_1,\ldots,H_{K_{\mathrm{comp}}}),\qquad
 H^{-1}=\blkdiag(H_1^{-1},\ldots,H_{K_{\mathrm{comp}}}^{-1}).
 \label{eq:block}
\end{equation}
Here $K_{\mathrm{comp}}$ counts connected components, while $K$ in
Section~\ref{sec:penalty} is a penalty matrix. The conclusion follows
from the factor $Z_{ir}Z_{ir'}$ and block multiplication.
Two columns can have a zero direct cross block while remaining in the same
component through a third column; that does not license separate inversion.
Cell-means indicators for disjoint groups separate, whereas a shared intercept
or overlapping continuous covariates generally connect the spatial columns.

The group-specific $\xi_g$ in Section~\ref{sec:extensions} likewise have no
cross-group spatial Hessian terms. Shared $\gamma$ coefficients usually
couple the groups through a border; the full covariance then need not be
block diagonal.

\subsection{Schur inversion when global coefficients form a border}
Let $d=GP$ be the number of grouped spatial coefficients; the number of
global coefficients is $R$. To keep the scalar intensity total $E_g$
distinct from an information block, write
\begin{equation}
 H=\begin{pmatrix}\mathcal D&\mathcal F\\
                 \mathcal F^\top&\mathcal E\end{pmatrix},\qquad
 \mathcal D=\blkdiag(D_1,\ldots,D_{K_{\mathrm{comp}}}),\quad
 \mathcal F\in\R^{d\times R},\qquad
 \mathcal S=\mathcal E-\mathcal F^\top\mathcal D^{-1}\mathcal F.
\end{equation}
If $\mathcal D$ and $\mathcal S$ are nonsingular,
\begin{equation}\label{eq:schur}
 H^{-1}=
 \begin{pmatrix}
 \mathcal D^{-1}+\mathcal D^{-1}\mathcal F\mathcal S^{-1}
                 \mathcal F^\top\mathcal D^{-1}
 &-\mathcal D^{-1}\mathcal F\mathcal S^{-1}\\
 -\mathcal S^{-1}\mathcal F^\top\mathcal D^{-1}&\mathcal S^{-1}
 \end{pmatrix}.
\end{equation}
Multiplying by $H$ proves the identity using
$\mathcal E=\mathcal S+\mathcal F^\top\mathcal D^{-1}\mathcal F$.
With $\mathcal F=0$ it reduces to independent
block inversion. Use solves with $D_k$ and $\mathcal S$ rather than explicitly
forming intermediate inverses whenever only contrasts or selected entries
are required.
Writing $d_k=\dim D_k$, factorisation costs
$O(\sum_kd_k^3+R^3)$; solving for $\mathcal D^{-1}\mathcal F$ costs
$O(\sum_kd_k^2R)$, forming the Schur complement costs $O(dR^2)$, and forming
the dense spatial covariance correction costs $O(d^2R)$.
An invertible $H$ can have a singular leading block $\mathcal D$, so invertibility of
$H$ alone is not sufficient for this particular Schur formula.

\subsection{Shared covariates induce cross-group covariance}\label{sec:shared}
For grouped CBMR take $D_g=H_{\xi_g\xi_g}$ and
$F_g=H_{\xi_g\gamma}$, so $\mathcal F$ stacks the $F_g$ vertically.
Under the observed-information covariance choice $V_\vartheta=H^{-1}$,
\begin{align}\label{eq:shared-covariance}
 (V_\vartheta)_{\gamma\gamma}&=\mathcal S^{-1},\\
 (V_\vartheta)_{\xi_g\xi_h}
 &=\mathbf1\{g=h\}D_g^{-1}
       +D_g^{-1}F_g\mathcal S^{-1}F_h^\top D_h^{-1}.
\end{align}
Thus a block-diagonal \emph{spatial information block} does not imply
independent spatial estimates after estimating shared $\gamma$.
For global-covariate inference, $\mathcal E^{-1}$ would treat the spatial
coefficients as known; $\mathcal S^{-1}$ accounts for their estimation.
The same block identities hold using $I$ instead of $H$, or for penalised
curvature after replacing $\mathcal D$ by its penalised version, but the
latter is not automatically the sampling covariance
\eqref{eq:penalty-covariance}.

This matters directly for Equation 14 of the published article: the small
dimension of $\gamma$ does not ensure nonsingular or well-conditioned
profile information. A global covariate confounded with group intercepts
can make $\mathcal S$ singular regardless of the magnitude of its
coefficient. Independently fitted groups have zero cross-covariance only
when their data and fitted parameter vectors are genuinely separate.

\subsection{Condition number from symmetric eigenvalues}
If $H=Q\Lambda Q^\top$ with orthogonal $Q$, then
$H^\top H=Q\Lambda^2Q^\top$. Consequently,
\begin{equation}\label{eq:condition}
 \kappa_2(H)=\frac{\max_a|\lambda_a(H)|}{\min_a|\lambda_a(H)|}.
\end{equation}
Positive definiteness is not needed for this identity. For a singular matrix
the finite inverse and finite condition number are unavailable.
For \eqref{eq:block}, the characteristic polynomial is the product of the
block characteristic polynomials, hence
\[
 \kappa_2(H)=
 \frac{\max_k\max_a|\lambda_a(H_k)|}
      {\min_k\min_a|\lambda_a(H_k)|}.
\]
This is not in general $\max_k\kappa_2(H_k)$.
For a bordered matrix, the eigenvalues of $\mathcal D$ and $\mathcal S$ are
\emph{not} the eigenvalues of $H$; use the full symmetric eigenproblem.
All these equalities are exact algebraically; numerical eigensolvers and
inverses still incur rounding error.

An appropriate computation uses structural blocks when available, otherwise
Schur inversion when its hypotheses hold, otherwise Cholesky for a
positive-definite matrix or a general factorisation for a nonsingular
indefinite matrix. The latter is an algebraic inverse, not evidence of a valid
positive covariance estimate. Failed factorisations or singularity must be
reported rather than replaced by an apparently valid covariance.

The source covariance note also considers LAPACK's \texttt{pocon}, which
estimates $\kappa_1$, not $\kappa_2$. Norm equivalence gives
$n^{-1}\kappa_1(H)\leq\kappa_2(H)\leq n\kappa_1(H)$ for an $n\times n$
matrix; it neither makes the norms equal nor bounds the estimator's error
by itself. A symmetric eigensolver targets the required 2-norm quantity.
The source notes' timing and memory measurements are empirical results,
not consequences of the symbolic proofs and not benchmarks of this appendix.

\section{Model-based and robust voxelwise inference}\label{sec:inference}

\subsection{SV-CBMR: voxelwise covariate-effect contrasts}
All quantities in this section are evaluated at the fitted parameters.
The lesion manuscript's model-based choice is
$V_\beta=\widehat{\Cov}(\widehat\beta)=I(\widehat\beta)^{-1}$.
An observed-information alternative uses $H(\widehat\beta)^{-1}$; they agree
for Poisson, not generally for NB.
For $\mathcal T_j=C\otimes B_j^\top$, the voxelwise log-contrast and covariance are
\begin{equation}\label{eq:contrast}
 \widehat\eta_{Cj}=\mathcal T_j\widehat\beta,\qquad
 V_{Cj}=\mathcal T_jV_\beta\mathcal T_j^\top.
\end{equation}
For $S$ linearly independent contrast rows with nonsingular $V_{Cj}$, the
quadratic statistic
$\widehat\eta_{Cj}^\top V_{Cj}^{-1}\widehat\eta_{Cj}$
has the usual asymptotic $\chi^2_S$ null law under the statistical regularity
conditions. When $S=1$, its signed square root is the Wald statistic.
The degrees of freedom must be reduced if the tested constraints are
redundant. A contrast need not have row sum one; difference contrasts
naturally have row sum zero.

\subsection{GC-CBMR: group-map contrasts and global covariate tests}
For the published CBMR predictor, let $C_G\in\R^{S\times G}$ contain group
contrasts and use the full covariance $V_\vartheta$ appropriate to the
fitting/inference method. Define
\begin{equation}\label{eq:group-contrast}
 \mathcal T_{Gj}=[\,C_G\otimes B_j^\top\ \ 0_{S\times R}\,],\qquad
 \widehat\delta_j=\mathcal T_{Gj}\widehat\vartheta-\delta_{0j},\qquad
 V_{\delta j}=\mathcal T_{Gj}V_\vartheta\mathcal T_{Gj}^\top.
\end{equation}
The reference value $\delta_{0j}$ specifies the null, such as a known
homogeneous log-intensity or a zero log difference. The statistic
$\widehat\delta_j^\top V_{\delta j}^{-1}\widehat\delta_j$ restates
Equation 12 of the article with dimensionally consistent notation.
For one group, a log-map covariance across voxels is
$B\,V_{\xi_g\xi_g}B^\top$; the voxel $j$ variance is
$B_j^\top V_{\xi_g\xi_g}B_j$, with no transpose reversal.

For the log-intensity difference between groups $g$ and $h$,
\begin{equation}
 \Var(\widehat\eta_{Bgj}-\widehat\eta_{Bhj})
 =B_j^\top
   (V_{\xi_g\xi_g}+V_{\xi_h\xi_h}
     -V_{\xi_g\xi_h}-V_{\xi_h\xi_g})B_j.
\end{equation}
The cross terms from Section~\ref{sec:shared} must be retained.
No balancing of $M_g$ and $M_h$ is required algebraically; unequal group
sizes enter the likelihood and its information.

If intensity rather than log-intensity is reported, the Jacobian for
$\mu_{Bg}=\exp(B\xi_g)$ is
$\diag(\mu_{Bg})B$. For the vector of group intensities at voxel $j$,
the delta-method covariance is
\[
 V_{\mu,j}\ \approx\
 \diag(\widehat\mu_{B1j},\ldots,\widehat\mu_{BGj})\,
 V_{\eta,j}\,
 \diag(\widehat\mu_{B1j},\ldots,\widehat\mu_{BGj}).
\]
Here $V_{\eta,j}=(I_G\otimes B_j^\top)
V_{\xi_{\mathrm{all}}\xi_{\mathrm{all}}}(I_G\otimes B_j)$.
The Jacobian identity is exact; the covariance transformation is a
first-order approximation. Log contrasts avoid that additional approximation.
For global contrasts, use $C_\gamma\widehat\gamma$ and
$C_\gamma V_{\gamma\gamma}C_\gamma^\top$, with
$V_{\gamma\gamma}=\mathcal S^{-1}$ under the observed-information choice,
not the inverse global diagonal block.

\subsection{Independent-unit sandwich covariance}
For independent subjects, allowing arbitrary within-subject dependence, let
$U_i$ be the family-specific score in \eqref{eq:nb-score} and define
\begin{equation}\label{eq:sandwich}
 \mathcal A=-\sum_i\frac{\partial U_i}{\partial\beta^\top},\qquad
 \mathcal J=\sum_i U_iU_i^\top,\qquad
 V_{\mathrm{robust}}=\mathcal A^{-1}\mathcal J\mathcal A^{-\top}.
\end{equation}
For likelihood scores the empirical sensitivity $\mathcal A$ is $H$.
Under a correctly specified conditional mean, the expected sensitivity is
$I$ for the independent NB working score, so using $I$ as bread is an
asymptotically justified alternative. For Poisson they coincide exactly.
For a general misspecified estimating equation, its own derivative must be
used; NB Fisher weights combined with unweighted Poisson scores are not
the matching likelihood sandwich.
These formulas treat dispersion as fixed. If dispersion is jointly estimated
and its uncertainty is to be propagated robustly, stack its estimating
equation with the regression score before forming bread and meat.

The subject-clustered choice of $\mathcal J$ is essential: replacing it by a
sum of separate voxel score outer products discards within-subject cross terms.
For a particular contrast one can avoid storing the full meat:
\begin{equation}
 \mathcal T_jV_{\mathrm{robust}}\mathcal T_j^\top
 =\sum_i v_{ij}v_{ij}^\top,\qquad
 v_{ij}=\mathcal T_j\mathcal A^{-1}U_i.
\end{equation}
This follows by distributing the matrix product over the subject sum.
The independent-cell scores above must not be reused unchanged for the
aggregated or clustered likelihoods of Section~\ref{sec:extensions}.
For a grouped penalised fit use its corresponding scores, covariance
\eqref{eq:penalty-covariance}, and projection $\mathcal T_{Gj}$.

\subsection{Identification and the limits of algebraic improvements}
In the separable model, if $Bc_B=\mathbf1_N$ and $Zc_Z=\mathbf1_M$,
then $(\xi,\gamma)\mapsto(\xi+t c_B,\gamma-t c_Z)$ leaves every mean
unchanged. A constraint or removal of a redundant intercept is therefore
needed before using \eqref{eq:schur}. Similarly, adding an intercept column
to a basis that already spans the constant function can create rank
deficiency. Neither a large sample size nor a well-behaved condition number
alone establishes valid asymptotic inference; identifiability, independent
clusters and adequate information in each tested direction are separate
requirements.

The article reports an empirical threshold of roughly 200 foci per group
for its parametric-inference experiments. This is not a universal theorem,
and exact Hessians do not establish a new validity threshold. Removing an
autodifferentiation intermediate or replacing an SVD changes computational
cost, not the information in a sparse dataset. Likewise, adding a small
diagonal term solely to make inversion succeed changes the inferential
object and is not a silent substitute for a specified penalty.

\subsection{Bootstrap remains a distinct inferential procedure}
The parametric-bootstrap approach in Section 2.1.4 of the published article
addresses calibration when its Wald approximation is unreliable. It is not
made unnecessary by a faster Hessian. A model-consistent bootstrap fits
the stated null, simulates under its observation/dependence model with the
original groups and covariates, refits with the same constraints and
penalty-selection rule, and recomputes the chosen statistic.
Shared gamma effects must be preserved for clustered NB; a matched
aggregate-NB sampling law and an independent-cell NB sampling law are not
interchangeable. This prescription is distinct from a conditional
spatial-homogeneity randomisation scheme and its assumptions.

Write $B_{\mathrm{boot}}$ for the number of replicates, reserving $B$ for the
spatial basis. A common nonzero Monte Carlo tail estimate is
\[
 \widehat p=\frac{1+\sum_{b=1}^{B_{\mathrm{boot}}}
       \mathbf1\{T_b^*\geq T_{\mathrm{obs}}\}}{B_{\mathrm{boot}}+1}.
\]
The tail and statistic must match the tested alternative. A fitted-null
parametric bootstrap is not thereby an exact finite-sample test.
The article's generalised-Pareto tail extrapolation introduces an additional
approximation that needs its own diagnostics. Neither bootstrap validity,
tail calibration nor the article's empirical threshold is established by
the symbolic checks here.

\section{Exact symbolic verification and reproducibility}\label{sec:verification}

The repository runs the four original CBMR proof modules together with
\texttt{proofs/appendix.py} and \texttt{proofs/published\_cbmr.py}.
Each claim is an expression difference reduced
by SymPy; acceptance requires the literal scalar zero or a zero matrix,
not a numerical tolerance. The generated report from the current run is
\begin{center}
% Generated by cbmr-proofs. Do not edit.
% Every expression below is rendered from the object sympy verified, not transcribed.
\begin{tabular}{lrr}
\toprule
Symbolic suite & Claims & Closed\\
\midrule
Poisson Hessian & 13 & yes\\
Negative Binomial Hessian & 27 & yes\\
Clustered Negative Binomial Hessian & 10 & yes\\
Covariance strategies & 9 & yes\\
Integrated appendix & 35 & yes\\
Published CBMR integration & 26 & yes\\
\midrule
Total & 120 & yes\\
\bottomrule
\end{tabular}

\end{center}

The original modules check the Poisson and clustered-NB coefficient
Hessians, the aggregated-NB reparameterisation and derivative components,
and the structural, spectral and inverse identities. The appendix adds
the lesion manuscript's independent NB score, observed and expected weights, binary
likelihood reduction, Poisson limits, coefficient ordering, generic
affine and separable Hessian assembly, Poisson Kronecker factorisation,
the distinction between constant and varying covariate effects and its
constant-function restriction, the NB nonseparability identity, clustered curvature, and projected
contrast/sandwich identities. The published-CBMR suite checks grouped
Poisson compression and conditional allocation, the NB matching moments,
clustered marginal factorisation and curvature, fixed quadratic penalties,
and shared-global covariance corrections. The generated report counts
algebraic claims; it is not a count of independently proved statistical
theorems.

\subsection{Where the published model is improved}
\begin{center}\small
\begin{tabularx}{\textwidth}{lTT}
\toprule
Published article & Integrated calculation & Boundary of the result\\
\midrule
Equations 1--5 & Group design \eqref{eq:group-design} and basis-product assembly &
Same predictor, no full stacked design allocation\\
Equation 6 & Both count marginals in \eqref{eq:group-poisson} &
Voxel totals alone omit the study-allocation score\\
Equations 7--9 & $\rho_g,A_g$ and \eqref{eq:aggregate-chain} &
Exact derivatives of a moment-matched likelihood\\
Equations 10--11 & \eqref{eq:cluster-factorisation} and
\eqref{eq:cluster-local-information} &
Exact marginalisation, including within-study dependence\\
Figure 2, Section 2.1.1 & Fixed penalty gradient, curvature and sandwich &
Smoothing bias and tuning uncertainty remain separate\\
Equations 12--14 & Group projections and Schur covariance &
Shared global effects couple spatial estimates\\
Section 2.1.4 & Conditioning diagnostics and bootstrap boundaries &
Faster algebra does not certify sparse-data Wald inference\\
\bottomrule
\end{tabularx}
\end{center}

\subsection{What is and is not machine checked}
Scalar identities hold for symbolic variables with the stated positivity
and reality assumptions. Explicit matrix checks use small fixed dimensions
with arbitrary symbolic entries: the appendix assembly uses
$M=N=P=R=2$, while the original likelihood checks use three voxels and two
subjects. Generic scalar functions check the affine chain-rule composition;
the aggregated-NB components are checked separately rather than simplifying
its entire composite Hessian in one pass. The original Schur check and the
appendix Kronecker inverse check use formal matrix symbols of arbitrary
conformable dimensions under invertibility assumptions.

The published-CBMR suite also uses small grouped and study-level examples
with symbolic counts and means (and a fixed allocation-score design).
Its penalty differentiation holds the roughness matrix and smoothing
weights fixed; the penalised covariance-difference identity is checked
in the scalar case.
The dimension-independent arguments are given in the text; fixed-size
symbolic instances are not proofs over all dimensions. The generated scalar
formulas and report are tied directly to the checked objects, but the
handwritten exposition is not mechanically certified. SymPy is a computer
algebra system, not a proof-assistant kernel. These checks also do not verify
NiMARE implementation code, floating-point stability, asymptotic
distributions, model adequacy, or empirical performance.

\subsection{Lean formalisation: a checked subset, not complete coverage}
The repository also contains a pinned Lean/Mathlib development in
\texttt{lean_proof_checks/}. It formalises finite-index predictor and stacked-design
identities, real-variable likelihood derivatives, vector quadratic-penalty
derivatives, NB algebra and structured covariance identities.
These proofs are checked by Lean rather than by symbolic
simplification in SymPy. Many finite-sum and matrix statements use arbitrary
finite index types instead of fixed matrix sizes.

\textbf{The full appendix has not yet been converted to Lean.}
The finite affine derivative-to-Hessian assembly, grouped Poisson and
clustered-NB curvature, and induced-norm spectral identities now have
formal proofs. The remaining aggregated-NB derivations, probabilistic
expectation and mixture results, and statistical inference still have
outstanding proof obligations. The section-by-section inventory is
\texttt{docs/lean\_coverage.md}. The SymPy table above must not be read as a
Lean theorem count.

\texttt{make lean} compiles the included proofs and audits their transitive
axiom dependencies. Admitted proofs and additional axioms are rejected;
the accepted foundations are \texttt{propext}, \texttt{Classical.choice}
and \texttt{Quot.sound}. Positivity, differentiability and invertibility
assumptions remain explicit in theorem statements. Unproved obligations
are recorded in the inventory, not asserted as axioms.

The equation-level mapping in \texttt{docs/lean\_coverage.json} covers
all 54 displayed mathematical blocks and additional prose derivations.
It records the referenced theorems and any missing steps.
\texttt{make lean-complete} additionally rejects partial or pending
coverage; a successful \texttt{make lean} alone does not establish completeness.
The audit output and coverage report are saved in
\texttt{.lake/lean-audit.log} and
\texttt{.lake/lean-coverage-report.json}.

\subsection{Rebuilding}
From the repository root:
\begin{verbatim}
pip install -r requirements.txt
make appendix
\end{verbatim}
This verifies all claims, regenerates the LaTeX fragments and compiles
\texttt{docs/cbmr\_appendix.pdf} twice using \texttt{pdflatex}.
Any failed claim stops the build. \texttt{make proofs} runs only the
symbolic checks. CI performs the same checks and rejects stale generated
fragments; it does not currently build PDFs.
The editable source is \texttt{docs/cbmr\_appendix.tex}; the reference PDFs
are left unchanged.

\clearpage
\phantomsection
\addcontentsline{toc}{section}{References}
\begin{thebibliography}{9}
\bibitem{lesion}
\emph{Scalable Spatial Model for Brain MRI Lesion Mask Data}.
Supplied working manuscript, \texttt{docs/cbmr\_paper.pdf},
Sections 2.2--2.4. Source of the principal notation in this appendix.

\bibitem{published}
Yu, Y., Hill-Bowen, L. D., Riedel, M. C., Bottenhorn, K., Laird, A. R.,
and Nichols, T. E. (2025).
\emph{CBMR: Coordinate-based meta-regression for group and covariate inference}.
Imaging Neuroscience, 3.
\href{https://doi.org/10.1162/IMAG.a.1057}{doi:10.1162/IMAG.a.1057}.
Supplied as \texttt{docs/imag.a.1057.pdf}; this appendix adapts the
mathematics rather than reproducing the article's exposition.

\bibitem{hessian}
\emph{Closed-Form Observed Information for Coordinate-Based Meta-Regression}.
Repository derivation note, \texttt{docs/cbmr\_hessian.tex}.

\bibitem{covariance}
\emph{From the Hessian to the Covariance: Generalising the CBMR standard-error step}.
Repository derivation note, \texttt{docs/cbmr\_covariance.tex}.
\end{thebibliography}

\end{document}
